\documentclass[10pt,a4paper]{amsart}
\usepackage[margin=19mm]{geometry}
\usepackage{amsmath,amssymb,mathtools}
\usepackage[scr=rsfs,cal=euler]{mathalfa}
\usepackage{xfrac}
\usepackage[unicode,hidelinks,bookmarks=false]{hyperref}
\newcommand{\ie}{i.e.\ }
\newcommand{\cf}{cf.\ }
\newcommand{\resp}{respectively\ }
\newcommand{\Subsection}{\subsection}
\providecommand{\nobreakdash}{\nobreak\hbox{-}}
\setlength{\emergencystretch}{2em}
\multlinegap0pt
\usepackage{amssymb}
\usepackage{bm}
\usepackage{tikz}
\usepackage{algorithm}\usepackage{algpseudocode}
\usepackage{xcolor}
\usepackage{csquotes}
\usepackage{dsfont}
\newtheorem{assumption}{Assumption}
\newtheorem{lemma}{Lemma}
\newcommand{\E}{\mathbb{E}}
\newcommand{\R}{\mathbb{R}}
\title{Bayesian Distributionally Robust Merton Problem with Nonlinear Wasserstein Projections}
\author{Jose Blanchet, Jiayi Cheng, Hao Liu,, Yang Liu}
\date{Source: arXiv:2512.01408v2 (CC BY 4.0)}
\begin{document}
\maketitle

\noindent\emph{Source and licence.} Bayesian Distributionally Robust Merton Problem with Nonlinear Wasserstein Projections. Version arXiv:2512.01408v2 (CC BY 4.0), distributed by arXiv under the Creative Commons Attribution 4.0 International licence, http://creativecommons.org/licenses/by/4.0/, as recorded in the arXiv licence metadata for this exact version and shown on the abstract page https://arxiv.org/abs/2512.01408v2. Full licence notice: \texttt{licenses/arxiv-2512.01408-CC-BY-4.0.txt}.\par\smallskip

\noindent\textit{Editorial scope.} Counted unit: the article's lemma lem:L22-fro (source0.tex lines 3403--3453). The article's complete original proof of it is delivered with it (source0.tex lines 3455--4392, one \texttt{proof} environment of 938 lines). No mathematics is written, repaired or re-derived: the statement and the proof are the article's own lines, in the article's own notation, and no retained line is substituted or re-flowed.  Retained uncounted as the source-local context the proof needs: lines 3281--3319; lines 3326--3336.  Nothing was omitted from any retained span, so the package carries no omission record.  No retained line contains a citation, so the wrapper carries no bibliography and no bibliography entry.  No retained line contains a figure environment, an image-inclusion command or drawing code, so no image is delivered and the package declares no asset dependency.  Editorial formatting only, in addition: an AMS \texttt{amsart} preamble that restates the article's own declarations and macros used by the retained lines, namely the environments \texttt{assumption}, \texttt{lemma} on separate counters, together with the standard packages those lines need.  The retained environments print in this document as Assumption 1, Lemma 1 in order of appearance, whereas the article prints the same items with the numbering of its own full document.  The complete native source of the article as distributed in the arXiv e-print for this exact version is bundled unchanged as \texttt{sources/190122/source0.tex}.
\begin{assumption}\label{subGaussianFinal}
Suppose there exists $\gamma_0>0$ such that
\[
\E_{\mathbb P_0}
\left[
\exp\!\left(\gamma^2\|B\|_2^2\right)
\right]
<\infty
\qquad
\text{for every }\gamma<\gamma_0,
\]
and
\[
\gamma_0^2
>
T\|\sigma^{-1}\|_F^2
\max\left\{
\frac12
\left(
\beta^2-\beta+2
+
\sqrt{
(\beta^2-\beta+2)^2+8(\beta^2+\beta)
}
\right),
\,8
\right\}.
\]

When $\beta>0$, assume additionally that there exist $r_\beta>\max\{1,2\beta\}$
and conjugate exponents $p_\beta,q_\beta>1$ satisfying $\frac1{p_\beta}+\frac1{q_\beta}=1.$
Assume further that $2T(r_\beta^2-r_\beta)
\|\sigma^{-1}\|_F^2\,p_\beta
<
\gamma_0^2$ and $2T(r_\beta^2-r_\beta)
\|\sigma^{-1}\|_F^2\,q_\beta
<
\tau.$
\end{assumption}

\begin{equation}\label{eq:Hb}
H(b)
:=
\frac{1}{1-\alpha}
\int_{\mathbb R^d}
\nabla_bL_T(b,y)
\left(
\E_{\mathbb P_0}[L_T(B,y)]
\right)^{\frac{\alpha}{1-\alpha}}
\varphi_T(y)\,dy
\end{equation}

\begin{lemma}\label{lem:L22-fro}
Under Assumption~\ref{subGaussianFinal},
\begin{itemize}
    \item[(1)]
    For any $\alpha<1$ and $\alpha\neq0$, the vector field in
    \eqref{eq:Hb} is well-defined $\mathbb P_0$-almost surely and
    satisfies $\E_{\mathbb P_0}
    \left[
    \|H(B)\|_2^2
    \right]
    <\infty.$

    \item[(2)]
    For any $s\in[0,1]$, let $(B,B+\Delta)$ be distributed according
    to a coupling $\pi$ whose first marginal is $\mathbb P_0$, and satisfying $\E_\pi
\left[
\exp\!\left(\tau\|\Delta\|_2^2\right)
\right]
\leq 2.$ We define
    \[
    F_s(B,\Delta)
    :=
    \int_{\mathbb R^d}
    \nabla_bL_T(B+s\Delta,y)
    \left(
    \E_{\pi}[L_T(B+s\Delta,y)]
    \right)^{\frac{\alpha}{1-\alpha}}
    \varphi_T(y)\,dy.
    \]

    Then there exist $\eta>0$, a constant $C_1>0$, and a measurable
    random variable $C(B)>0$ such that
    \[
    2(1+\eta)C_1<\tau,
    \qquad
    \E_{\mathbb P_0}
    \left[
    C(B)^{2(1+\eta)}
    \right]
    <\infty,
    \]
    and, $\pi$-almost surely,
    \[
    \|F_s(B,\Delta)\|_2^2
    \leq
    C(B)
    \exp\!\left(C_1\|\Delta\|_2^2\right),
    \qquad s\in[0,1].
    \]
\end{itemize}
\end{lemma}

\begin{proof}
We first prove part~(1). Recall that $\beta:=\frac{\alpha}{1-\alpha}.$
Define $A(y):=\E_{\mathbb P_0}[L_T(B,y)].$
We begin with the case $\alpha<0$.
Direct computations give
\[
\nabla_b L_T(b,y)
=
L_T(b,y)\,\sigma^{-T}\!\bigl(y-Ta(b)\bigr),
\]
and
\[
L_T(b,y)^2\varphi_T(y)
=
\exp\!\bigl(T\|a(b)\|_2^2\bigr)
\varphi_T\!\bigl(y-2Ta(b)\bigr).
\]
By the Jensen inequality,
\[
A(y) \;=\; \E_{\mathbb{P}_0}\big[e^{\langle a(B),y\rangle-\frac{T}{2}\|a(B)\|_2^2}\big]
\;\ge\;
\exp\!\Big( \langle \E_{\mathbb{P}_0}[a(B)],y\rangle - \tfrac{T}{2}\,\E_{\mathbb{P}_0}\left[\|a(B)\|_2^2\right] \Big).
\]
Raising to the negative power $\beta<0$ reverses the inequality and gives
\begin{equation}\label{eq:est_A_beta}
A(y)^{\beta}\;\le\;C_a\,\exp\!\big(\beta\,\langle v,y\rangle\big),
\qquad
C_a:=\exp\!\Big(-\tfrac{\beta T}{2}\,\E_{\mathbb{P}_0}\left[\|a(B)\|_2^2\right]\Big),\ \ 
v:=\E_{\mathbb{P}_0}[a(B)].
\end{equation}
Since
\[
\|\nabla_b L_T(B,y)\|_2^2
= L_T(B,y)^2\,\|\sigma^{-T}(y-T a(B))\|_{2}^2
\le \,\|\sigma^{-1}\|_F^2\,L_T(B,y)^2\,\|y-T a(B)\|_2^2,
\] 
we use $\|u+v\|_2^2\le 2(\|u\|_2^2+\|v\|_2^2)$ and have %
\begin{equation*}
\|\nabla_b L_T(B,y)\|_2^2
\;\le\; C\,\|\sigma^{-1}\|_F^2\,L_T(B,y)^2\,\big(1+\|y\|_2^2+T^2\|a(B)\|_2^2\big),
\end{equation*}where $C > 0$ is a constant (we use $C$ to absorb constants if there is no confusion). Therefore,
\[
\begin{aligned}
&\int_{\R^d} A(y)^{\beta}\,\E_{\mathbb{P}_0}\big[\|\nabla_b L_T(B,y)\|_2^2\big]\;\varphi_T(y)\,dy \\
&\ \ \le\ C\,\|\sigma^{-1}\|_F^2\,\E_{\mathbb{P}_0}\!\Big[\int_{\R^d} L_T(B,y)^2\,(1+\|y\|_2^2+T^2\|a(B)\|_2^2)\,e^{\beta\langle v,y\rangle}\,\varphi_T(y)\,dy\Big]\\
&\ \ =\ C\,\|\sigma^{-1}\|_F^2\,\E_{\mathbb{P}_0}\!\Big[e^{T\|a(B)\|_2^2}
\int_{\R^d} (1+\|y\|_2^2+T^2\|a(B)\|_2^2)\,e^{\beta\langle v,y\rangle}\,\varphi_T\!\big(y-2T a(B)\big)\,dy\Big].
\end{aligned}
\]
The inner integral is 
\begin{align*}
    &\int_{\mathbb{R}^d} \bigl(1+\|y\|_2^{2}+T^{2}\|a(B)\|_2^{2}\bigr)\,
e^{\beta\langle v,y\rangle}\,
\varphi_{T}\bigl(y-2T\,a(B)\bigr)\,dy\\
&=
\exp\!\Big(2T\beta\,\langle v,a(B)\rangle + \tfrac{T}{2}\|\beta v\|_2^{2}\Big)\,
\Big[\,1 + T d + \|2T a(B) + T\beta v\|_2^{2} + T^{2}\|a(B)\|_2^{2}\,\Big].
\end{align*}
For a fixed $\varepsilon > 0$ (which will be chosen later), by the Young inequality, 
\[
2T\beta\,\langle v,a(B)\rangle
\leq
2T|\beta|\,|\langle v,a(B)\rangle|
\leq
\varepsilon\|a(B)\|_2^2
+
\frac{T^2\beta^2\|v\|_2^2}{\varepsilon}.
\]
Hence, 
\[
e^{\,2T\beta\langle v,a(B)\rangle+\frac{T}{2}\|\beta v\|_2^2}
\le C_{\mathrm{tilt}}(\varepsilon)\,e^{\,\varepsilon\|a(B)\|_2^2},
\]
with
\[
C_{\mathrm{tilt}}(\varepsilon):=\exp\!\Big(\tfrac{T}{2}\|\beta v\|_2^2+\tfrac{T^2\beta^2\|v\|_2^2}{\varepsilon}\Big).
\]
Using the Young inequality in the quadratic bracket, we have
\[
\|2Ta+T\beta v\|_2^2\le (1+\varepsilon)\,(2T)^2\|a\|_2^2+\Big(1+\tfrac1\varepsilon\Big)T^2\beta^2\|v\|_2^2,
\]
and further get
\[
1+Td+\|2Ta(B)+T\beta v\|_2^2+T^2\|a(B)\|_2^2
\le C_0(\varepsilon)+C_1(\varepsilon)\,\|a(B)\|_2^2,
\]where $C_0(\varepsilon):=1+Td+\Big(1+\tfrac{1}{\varepsilon}\Big)T^2\beta^2\|v\|_2^2,$ and $C_1(\varepsilon):=T^2\big(5+4\varepsilon\big).$
Therefore, the inner integral is upper bounded by
$C_{\mathrm{tilt}}(\varepsilon)\,\bigl(C_0(\varepsilon)+C_1(\varepsilon)\,\|a(B)\|_2^2\bigr)\,e^{\,\varepsilon\|a(B)\|_2^2}.$

Therefore, there exists a constant $C > 0$ such that
\begin{align*}
    &\int_{\R^d} A(y)^{\beta}\,\E_{\mathbb{P}_0}\big[\|\nabla_b L_T(B,y)\|_2^2\big]\;\varphi_T(y)\,dy \\
    &\leq C\,\|\sigma^{-1}\|_F^2\,C_{\mathrm{tilt}}(\varepsilon)\;
\E_{\mathbb{P}_0}\!\Big[\bigl(C_0(\varepsilon)+C_1(\varepsilon)\|a(B)\|_2^2\bigr)\,e^{\,(T+\varepsilon)\|a(B)\|_2^2}\Big].
\end{align*}

Let
\[
\varepsilon_*:=\tfrac12\!\left(\frac{\gamma_0^2}{2\|\sigma^{-1}\|_F^2}-T\right)
\quad\text{(well-defined and $>0$ if }\gamma_0^2>2T\|\sigma^{-1}\|_F^2).
\]
Since $a(B)=\sigma^{-1}B-m$, we have $\|a(B)\|_2^2\le 2\|\sigma^{-1}\|_F^2\|B\|_2^2+2\|m\|_2^2$. Hence, 
\begin{equation}\label{goal}
\bigl(C_0(\varepsilon_*)+C_1(\varepsilon_*)\|a(B)\|_2^2\bigr)\,e^{\,(T+\varepsilon_*)\|a(B)\|_2^2}
\;\le\;K_m\,(1+c\|B\|_2^2)\,e^{\,C'\|B\|_2^2},    
\end{equation}
for constants $K_m,c>0$ and
\[
C':=2\,(T+\varepsilon_*)\,\|\sigma^{-1}\|_F^2
= T\|\sigma^{-1}\|_F^2+\frac{\gamma_0^2}{2}\;<\gamma_0^2.
\]

To see (\ref{goal}), recall that \(a(B)=\sigma^{-1}B - m\), and denote
\[
\kappa_1 := 2\,\|\sigma^{-1}\|_F^2,\qquad \kappa_0 := 2\,\|m\|_2^2.
\]
Then
$$\|a(B)\|_2^2 \;=\; \|\sigma^{-1}B - m\|_2^2
\;\le\; 2\|\sigma^{-1}B\|_2^2 + 2\|m\|_2^2
\;\le\; \kappa_1 \|B\|_2^2 + \kappa_0,$$
and
\begin{align*}
(T+\varepsilon_*)\|a(B)\|_2^2
&\le (T+\varepsilon_*)\big(\kappa_1 \|B\|_2^2 + \kappa_0\big)
= \underbrace{2(T+\varepsilon_*)\|\sigma^{-1}\|_F^2}_{=:~C'} \|B\|_2^2 \;+\; 2(T+\varepsilon_*)\|m\|_2^2,
\end{align*}
where
$C'= 2(T+\varepsilon_*)\|\sigma^{-1}\|_F^2
= T\|\sigma^{-1}\|_F^2 + \frac{\gamma_0^2}{2} \;<\; \gamma_0^2.$ 
Hence, 
$$e^{(T+\varepsilon_*)\|a(B)\|_2^2}
\;\le\; e^{\,2(T+\varepsilon_*)\|m\|_2^2}\; e^{\,C'\|B\|_2^2}.$$
For the polynomial prefactor, 
\begin{align*}
C_0(\varepsilon_*) + C_1(\varepsilon_*)\|a(B)\|_2^2
&\le C_0(\varepsilon_*) + C_1(\varepsilon_*)\big(\kappa_1 \|B\|_2^2 + \kappa_0\big) \\
&= \underbrace{\big(C_0(\varepsilon_*) + C_1(\varepsilon_*)\kappa_0\big)}_{=:~K_1}
\;+\; \underbrace{C_1(\varepsilon_*)\kappa_1}_{=:~K_2}\,\|B\|_2^2 \\
&= K_1\Big(1 + \frac{K_2}{K_1}\,\|B\|_2^2\Big)
= K_1\big(1 + c\,\|B\|_2^2\big), 
\end{align*}
where
\[
c := \frac{C_1(\varepsilon_*)\kappa_1}{C_0(\varepsilon_*) + C_1(\varepsilon_*)\kappa_0} \;>\; 0 .
\]
Therefore, we get \eqref{goal} with
\[
K_m \;:=\; K_1\, e^{\,2(T+\varepsilon_*)\|m\|_2^2}
\;=\; \big(C_0(\varepsilon_*) + C_1(\varepsilon_*)\kappa_0\big)\, e^{\,2(T+\varepsilon_*)\|m\|_2^2}.
\]
Since \(C' < \gamma_0^2\), we pick any \(\gamma\) such that \(C' < \gamma^2 < \gamma_0^2\).
For \(t\ge 0\) and \(\delta := \gamma^2 - C' > 0\), the elementary bound
\[
1 + c t \;\le\; \Big(1 + \frac{c}{e\,\delta}\Big)\, e^{\delta t}
\]
implies
\[
(1 + c\,\|B\|_2^2)\, e^{\,C'\|B\|_2^2}
\;\le\; K\, e^{\,\gamma^2 \|B\|_2^2},
\qquad K := 1 + \frac{c}{e(\gamma^2 - C')}.
\]
Therefore, under the sub-Gaussian assumption
\(\E_{\mathbb{P}_0}\big[e^{\gamma^2\|B\|_2^2}\big] < \infty\) for all \(\gamma < \gamma_0\),
\[
\E_{\mathbb{P}_0}\!\Big[\bigl(C_0(\varepsilon_*) + C_1(\varepsilon_*)\|a(B)\|_2^2\bigr)\, e^{(T+\varepsilon_*)\|a(B)\|_2^2}\Big]
\;\le\; K_m\,K\,\E_{\mathbb{P}_0}\big[e^{\gamma^2\|B\|_2^2}\big] \;<\; \infty.
\]
In order to show the well-definedness of $H$, we need to show
\begin{equation}\label{CS}
    \int_{\mathbb{R}^d}A(y)^{\beta}\varphi_T(y)dy < \infty
\end{equation}since from the Cauchy-Schwarz inequality, for $\mathbb{P}_0$ almost every $b$,
\begin{align*}
    &\left\Vert H(b) \right\Vert_2 \leq \frac{1}{1-\alpha}\int_{\mathbb R^d}A(y)^{\beta}\left\Vert\nabla_b L_T(b,y)\right\Vert_2\,\,\varphi_T(y)\,dy\\
    &\leq \frac{1}{1-\alpha} \left(\int_{\mathbb{R}^d}\left\Vert\nabla_b L_T(b,y)\right\Vert_2^2A(y)^{\beta}\,\varphi_T(y)dy\right)^{\frac{1}{2}}\left(\int_{\mathbb{R}^d}A(y)^{\beta}\,\varphi_T(y)dy\right)^{\frac{1}{2}},
\end{align*}
{The first term is finite $\mathbb P_0$-almost surely. Squaring, taking
$\mathbb P_0$-expectations, and applying Tonelli's theorem complete the proof.}

When $\alpha < 0$, we recall Eq. \eqref{eq:est_A_beta} and have %
\[
\int_{\R^d} A(y)^{\beta}\,\varphi_T(y)\,dy
\;\le\;
C_a \int_{\R^d} e^{\beta\langle v,y\rangle}\,\varphi_T(y)\,dy.
\]
Completion of squares gives
\[
\int_{\R^d} e^{\beta\langle v,y\rangle}\,\varphi_T(y)\,dy
=(2\pi T)^{-d/2}\!\!\int_{\R^d}
\exp\!\left(\beta\langle v,y\rangle-\frac{\|y\|_2^2}{2T}\right)dy
=\exp\!\left(\frac{T}{2}\|\beta v\|_2^2\right).
\]
Indeed,
\[
\beta\langle v,y\rangle-\frac{\|y\|_2^2}{2T}
=-\frac{1}{2T}\Big\|\,y-T\beta v\,\Big\|_2^2+\frac{T}{2}\|\beta v\|_2^2.
\]
Hence, 
\[
\int_{\R^d} A(y)^{\beta}\,\varphi_T(y)\,dy
\;\le\;
C_a\,\exp\!\left(\frac{T}{2}\|\beta v\|_2^2\right)
<\infty.
\]
This finishes the proof of the case when $\alpha<0$.
Next we divide the positive $\alpha$ into two cases.

\noindent\emph{Case 1: $0<\beta\le 1$ (equivalently $0<\alpha\le \tfrac12$).}
Let $\lambda_T(dy):=\varphi_T(y)\,dy$; then $\lambda_T$ is a probability measure since
\[
\int_{\R^d}\varphi_T(y)\,dy=1.
\]
Set $f(y):=A(y)=\E_{\mathbb P_0}[L_T(B,y)]\geq0$ and
\[
\phi(x):=x^{\beta},\qquad x\ge 0.
\]
For $0<\beta\le 1$, the map $\phi$ is concave and increasing, so the Jensen 
inequality for concave functions gives
\[
\phi\!\left(\int_{\R^d} f(y)\,\lambda_T(dy)\right)\;\ge\;\int_{\R^d}\phi(f(y))\,\lambda_T(dy),
\]
i.e.
\begin{equation}\label{eq:jensen-easy}
\int_{\R^d} A(y)^{\beta}\,\varphi_T(y)\,dy
\;\le\;
\Big(\int_{\R^d} A(y)\,\varphi_T(y)\,dy\Big)^{\beta}.
\end{equation}

It remains to compute $\int_{\R^d} A(y)\,\varphi_T(y)\,dy$. By the Fubini theorem,
\[
\int_{\R^d} A(y)\,\varphi_T(y)\,dy
=\E_{\mathbb{P}_0}\!\left[\int_{\R^d} L_T(B,y)\,\varphi_T(y)\,dy\right] = 1
\]
from a completion of squares argument. Thus, plugging into \eqref{eq:jensen-easy} yields the bound
\[\ \int_{\R^d} A(y)^{\beta}\,\varphi_T(y)\,dy \;\le\; 1,\qquad 0<\beta\le 1. \ 
\]

\medskip
\noindent\emph{Case 2: $\beta>1$ (equivalently $\tfrac12<\alpha<1$).}
Since the map $x\mapsto x^{\beta}$ is convex on $(0,\infty)$, the Jensen inequality yields
\[
A(y)^{\beta}=\big(\E_{\mathbb{P}_0}[L_T(B,y)]\big)^{\beta}\le \E_{\mathbb{P}_0}\big[L_T(B,y)^{\beta}\big].
\]
Integrating and exchanging expectation and integral,
\[
\int_{\R^d} A(y)^{\beta}\,\varphi_T(y)\,dy
\;\le\;
\E_{\mathbb{P}_0}\!\left[\int_{\R^d} L_T(B,y)^{\beta}\,\varphi_T(y)\,dy\right].
\]
For fixed $b$, completing the square gives
\[
\int_{\R^d} L_T(b,y)^{\beta}\,\varphi_T(y)\,dy
=\exp\!\left(\frac{T}{2}(\beta^2-\beta)\,\|a(b)\|_2^2\right).
\]
Therefore
\[
\int_{\R^d} A(y)^{\beta}\,\varphi_T(y)\,dy
\;\le\;
\E_{\mathbb{P}_0}\!\left[\exp\!\left(\frac{T}{2}(\beta^2-\beta)\,\|a(B)\|_2^2\right)\right].
\]
Using $\|a(B)\|_2^2\le 2\|\sigma^{-1}\|_F^2\,\|B\|_2^2+2\|m\|_2^2$, we obtain
\begin{equation}\label{eq:est_int_A_beta}
\int_{\R^d} A(y)^{\beta}\,\varphi_T(y)\,dy
\;\le\;
\exp\!\Big(T(\beta^2-\beta)\,\|m\|_2^2\Big)\;
\E_{\mathbb{P}_0}\!\left[\exp\!\Big(T(\beta^2-\beta)\,\|\sigma^{-1}\|_F^2\,\|B\|_2^2\Big)\right].
\end{equation}
Hence, under the sub-Gaussian moment assumption
\[
\E_{\mathbb{P}_0}\!\left[e^{\gamma^2\|B\|_2^2}\right]<\infty\quad\text{for all }\gamma<\gamma_0,
\]
the right-hand side of Eq. \eqref{eq:est_int_A_beta} is finite whenever
\begin{equation}\label{eq:easy-threshold}
\ T(\beta^2-\beta)\,\|\sigma^{-1}\|_F^2 \;<\; \gamma_0^2.
\end{equation}

Next, we focus on the bound when $\alpha\in(0,1)$. Define
\[
\mathcal I_\beta
:=
\int_{\R^d}
A(y)^\beta
\E_{\mathbb P_0}
\left[
\|\nabla_bL_T(B,y)\|_2^2
\right]
\varphi_T(y)\,dy.
\]
Write
\[
G(y)
:=
\E_{\mathbb P_0}
\left[
\|\nabla_bL_T(B,y)\|_2^2
\right],
\]
and set
\[
I_0
:=
\int_{\R^d}G(y)\varphi_T(y)\,dy,
\qquad
I_1
:=
\int_{\R^d}A(y)G(y)\varphi_T(y)\,dy.
\]

When $0<\beta<1$, H\"older's inequality with exponents
$p=1/\beta$ and $q=1/(1-\beta)$ gives
\begin{equation}\label{eq:Ibeta-interpolate}
\begin{aligned}
\mathcal I_\beta
&=
\int_{\R^d}
\bigl(A(y)G(y)\bigr)^\beta
G(y)^{1-\beta}
\varphi_T(y)\,dy \\
&\leq
\left(
\int_{\R^d}A(y)G(y)\varphi_T(y)\,dy
\right)^\beta
\left(
\int_{\R^d}G(y)\varphi_T(y)\,dy
\right)^{1-\beta} \\
&=
I_1^\beta I_0^{1-\beta}.
\end{aligned}
\end{equation}
When $\beta=1$, we simply have
\[
\mathcal I_1=I_1.
\]
Thus, it suffices to show that
\[
I_0<\infty
\qquad\text{and}\qquad
I_1<\infty.
\]

Recall that there exists a constant $C > 0$ such that
\[
\|\nabla_b L_T(B,y)\|_2^2
\;\le\;C\,\|\sigma^{-1}\|_F^2\,L_T(B,y)^2\,(1+\|y\|_2^2+T^2\|a(B)\|_2^2).
\]
Using the Tonelli theorem,
\[
\begin{aligned}
I_0
&\le C\,\|\sigma^{-1}\|_F^2\,\E_{\mathbb{P}_0}\!\Big[e^{T\|a(B)\|_2^2}
\!\!\int_{\R^d}\!(1+\|y\|_2^2+T^2\|a(B)\|_2^2)\,\varphi_T(y-2Ta(B))\,dy\Big].
\end{aligned}
\]For the inner integral, 
fix \(a\in\R^d\) and set the notation
\[
J(a):=\int_{\R^d}\bigl(1+\|y\|_2^2+T^2\|a\|_2^2\bigr)\,\varphi_T(y-2Ta)\,dy.
\]
Let \(Z\sim\mathcal N(0,T I_d)\). Since \(\varphi_T(y-2Ta)\,dy\) is the law of \(Y:=Z+2Ta\), we have 
\[
J(a)=\E\!\left[1+\|Y\|_2^2+T^2\|a\|_2^2\right]
=1+\E\left[\|Z+2Ta\|_2^2\right]+T^2\|a\|_2^2 = 1+dT+5T^2\|a\|_2^2.
\]
With \(a=a(B)\), the inner integral in \(I_0\) equals \(1+dT+5T^2\|a(B)\|_2^2\), and therefore
\[
I_0 \le C\,\|\sigma^{-1}\|_F^2\,
\E_{\mathbb{P}_0}\!\Big[e^{T\|a(B)\|_2^2}\big(1+dT+5T^2\|a(B)\|_2^2\big)\Big].
\]
The estimate for \(I_0\) requires
\[
\mathbb{E}_{\mathbb{P}_0}\!\left[\exp\!\big(T\|a(B)\|_2^{2}\big)\right] < \infty,
\]
which is actually a sub-Gaussian assumption on \(B\) with condition
\[
2T\,\|\sigma^{-1}\|_F^{2} \;<\; \gamma_0^{2}.
\]
Next, we upper bound $I_1$. By the Young inequality \(\langle u,v\rangle\le \tfrac{\varepsilon}{2}\|u\|_2^2+\tfrac{1}{2\varepsilon}\|v\|_2^2\) with a fixed $\varepsilon > 0$ (which will be chosen later),
\[
A(y)\;\le\;\E_{\mathbb{P}_0}\!\Big[e^{-\frac{T-\varepsilon}{2}\|a(B)\|_2^2}\Big]\;
\exp\!\Big(\frac{\|y\|_2^2}{2\varepsilon}\Big)
\;=\;C_A(\varepsilon)\,e^{c\|y\|_2^2},
\qquad c:=\frac{1}{2\varepsilon}\in\Big(0,\frac{1}{2T}\Big).
\]
Here \(C_A(\varepsilon)=\E_{\mathbb{P}_0}\!\big[e^{\frac{\varepsilon-T}{2}\|a(B)\|_2^2}\big]\).
Hence, by the Tonelli theorem,
\[
\begin{aligned}
I_1
&\le C\,\|\sigma^{-1}\|_F^2\,C_A(\varepsilon)\,
\E_{\mathbb{P}_0}\!\Big[e^{T\|a(B)\|_2^2}\!\int_{\R^d}(1+\|y\|_2^2+T^2\|a(B)\|_2^2)e^{c\|y\|_2^2}\varphi_T(y-2Ta(B))\,dy\Big].
\end{aligned}
\]
Fix \(a\in\R^d\) and \(c\in(0,\frac{1}{2T})\). Let
\[
K_c(a):=\int_{\R^d}\bigl(1+\|y\|_2^2+T^2\|a\|_2^2\bigr)\,e^{c\|y\|_2^2}\,\varphi_T(y-2Ta)\,dy.
\]
With \(Y\sim\mathcal N(2Ta,TI_d)\) we have \(K_c(a)=\E[(1+\|Y\|_2^2+T^2\|a\|_2^2)e^{c\|Y\|_2^2}]\).
Define
\[
F(c,a):=\int_{\R^d} e^{c\|y\|_2^2}\,\varphi_T(y-2Ta)\,dy
=(1-2cT)^{-d/2}\exp\!\Big(\frac{4T^2 c}{1-2cT}\,\|a\|_2^2\Big).
\]
Then
\[
\frac{\partial}{\partial c}F(c,a)
=\int_{\R^d}\|y\|_2^2 e^{c\|y\|_2^2}\varphi_T(y-2Ta)\,dy
=F(c,a)\!\left(\frac{dT}{1-2cT}+\frac{4T^2\|a\|_2^2}{(1-2cT)^2}\right),
\]
so
\[
\begin{aligned}
K_c(a)
&=(1+T^2\|a\|_2^2)\,F(c,a)+\frac{\partial}{\partial c}F(c,a)\\
&=F(c,a)\!\left[1+T^2\|a\|_2^2+\frac{dT}{1-2cT}+\frac{4T^2\|a\|_2^2}{(1-2cT)^2}\right]
\;\le\;C_1(c,T,d)\,\bigl(1+\|a\|_2^2\bigr)\,F(c,a).
\end{aligned}
\]
Therefore, with \(c=\frac{1}{2\varepsilon}\in(0,\frac{1}{2T})\),
\[
\begin{aligned}
I_1
&\le C\,\|\sigma^{-1}\|_F^2\,C_A(\varepsilon)\,
\E_{\mathbb{P}_0}\!\Big[e^{T\|a(B)\|_2^2}\,K_c\big(a(B)\big)\Big]\\
&\le C(\varepsilon,T,d,\sigma)\,
\E_{\mathbb{P}_0}\!\Big[(1+\|a(B)\|_2^2)\,
\exp\!\Big(\Big(T+\frac{2T^2}{\varepsilon-T}\Big)\|a(B)\|_2^2\Big)\Big],
\end{aligned}
\]where \[
C(\varepsilon,T,d,\sigma)
= C\;\|\sigma^{-1}\|_{F}^{2}\;C_{A}(\varepsilon)\;
\Big(1-\tfrac{T}{\varepsilon}\Big)^{-d/2}
\left(1+\frac{dT}{1-\tfrac{T}{\varepsilon}}+T^{2}
+\frac{4T^{2}}{\big(1-\tfrac{T}{\varepsilon}\big)^{2}}\right).
\]
To make the bound for \(I_1\) finite, we only need an \(\varepsilon>T\) such that 
\[
\frac{\varepsilon-T}{2}\;<\;\frac{\gamma_0^2}{2\|\sigma^{-1}\|_F^2}
\qquad\text{and}\qquad
T+\frac{2T^2}{\varepsilon-T}\;<\;\frac{\gamma_0^2}{2\|\sigma^{-1}\|_F^2}.
\]
These two inequalities simultaneously hold whenever
\[
4T^{2}\,\|\sigma^{-1}\|_{F}^{4}\;+\;2T\,\|\sigma^{-1}\|_{F}^{2}\,\gamma_{0}^{2}\;<\;\gamma_{0}^{4},
\]which is equivalent to 
\[
 \gamma_0^2 \;>\; (1+\sqrt{5})\,T\,\|\sigma^{-1}\|_F^2.
\]
Under this %
condition%
, an explicit admissible choice is
\[
\varepsilon^{*}\;:=\;T+\frac{1}{2}\left(
\frac{2T^{2}}{\frac{\gamma_{0}^{2}}{2\|\sigma^{-1}\|_{F}^{2}}-T}
+\frac{\gamma_{0}^{2}}{\|\sigma^{-1}\|_{F}^{2}}
\right)\;>\;T.
\]
For this \(\varepsilon^{*}\) we have
\[
C_{A}(\varepsilon^{*})
=\mathbb{E}_{\mathbb{P}_{0}}\!\left[\exp\!\left(\frac{\varepsilon^{*}-T}{2}\,\|a(B)\|_{2}^{2}\right)\right]<\infty,
\qquad
\mathbb{E}_{\mathbb{P}_0}\!\left[(1+\|a(B)\|_2^2)\,
\exp\!\left(\Big(T+\frac{2T^2}{\varepsilon^{*}-T}\Big)\|a(B)\|_2^2\right)\right]<\infty,
\]
by the sub-Gaussian integrability above. Hence,  \(I_1<\infty\).

Finally, we upper bound $\mathcal I_\beta$ when $\beta > 1$.
By convexity of the map $x \mapsto x^{\beta}$ when $\beta > 1$ and the Jensen inequality, 
\[
A(y)^{\beta}\;=\;\big(\mathbb E_{\mathbb P_0}[L_T(B,y)]\big)^{\beta}\;\le\;\mathbb E_{\mathbb P_0}\!\big[L_T(B,y)^{\beta}\big].
\]
Fix \(\varepsilon>\beta^2 T\) and use the Young inequality
\(\langle u,v\rangle\le \tfrac{\varepsilon}{2}\|u\|_2^2+\tfrac{1}{2\varepsilon}\|v\|_2^2\) with \(u=\beta a(B)\), \(v=y\):
\[
L_T(B,y)^{\beta}
=\exp\!\Big(\beta\langle a(B),y\rangle-\tfrac{\beta T}{2}\|a(B)\|_2^2\Big)
\le \exp\!\Big(\tfrac{\varepsilon-\beta T}{2}\|a(B)\|_2^2\Big)\,e^{c\|y\|_2^2},
\quad c:=\frac{\beta^2}{2\varepsilon}\in\Big(0,\frac{1}{2T}\Big).
\]
Taking expectation in \(B\) yields
\[
A(y)^{\beta}\;\le\;C_{\beta}(\varepsilon)\,e^{c\|y\|_2^2},
\quad
C_{\beta}(\varepsilon):=
\mathbb E_{\mathbb P_0}\!\Big[\exp\!\Big(\tfrac{\varepsilon-\beta T}{2}\|a(B)\|_2^2\Big)\Big].
\]
Thus, exactly as in the \(I_1\) bound for \(\beta\in(0,1)\), we obtain 
\[
\mathcal I_\beta
\;\le\;C(\varepsilon,T,d,\sigma)\;
\mathbb E_{\mathbb P_0}\!\Big[(1+\|a(B)\|_2^2)\,
\exp\!\Big(\Big(T+\frac{2T^2\beta^2}{\varepsilon-\beta^2T}\Big)\|a(B)\|_2^2\Big)\Big],
\]
for some finite constant \(C(\varepsilon,T,d,\sigma)\) whenever \(c\in(0,1/(2T))\).
Therefore, to have \(\mathcal I_\beta<\infty\), it suffices to pick \(\varepsilon>\beta^2T\) such that
\[
\frac{\varepsilon-\beta T}{2}\;<\;\frac{\gamma_0^2}{2\|\sigma^{-1}\|_F^2}
\qquad\text{and}\qquad
T+\frac{2T^2\beta^2}{\varepsilon-\beta^2T}\;<\;\frac{\gamma_0^2}{2\|\sigma^{-1}\|_F^2}.
\tag{$\star$}
\]
These two inequalities simultaneously hold precisely when
\[
\frac{\gamma_0^2}{\|\sigma^{-1}\|_F^2}\;>\;\frac{T}{2}\Big(\beta^2-\beta+2
\;+\;\sqrt{\,(\beta^2-\beta+2)^{2}+8(\beta^2+\beta)\,}\Big)\;.
\]
Under the displayed condition, choose any
\[
\varepsilon\in\Big(\,\beta^2T+\frac{2T^2\beta^2}{\frac{\gamma_0^2}{2\|\sigma^{-1}\|_F^2}-T}\,,\;\beta T+\frac{\gamma_0^2}{\|\sigma^{-1}\|_F^2}\,\Big),
\]
which is a nonempty interval; then both parts of \((\star)\) hold and the expectation above is finite. Hence \(\mathcal I_\beta<\infty\) for all \(\beta>1\) under the stated sub-Gaussian slack.

For the proof of part~(2), let $\pi$ be a coupling whose first marginal
is $\mathbb P_0$, write $(B,B+\Delta)\sim\pi,$
and suppose that $\E_\pi
\left[
\exp\!\left(\tau\|\Delta\|_2^2\right)
\right]
\leq 2.$
Throughout this part, expectations involving both $B$ and $\Delta$ are
taken with respect to $\pi$.

For $s\in[0,1]$, define
\[
I_\beta^\pi(s)
:=
\int_{\R^d}
\left(
\E_\pi[L_T(B+s\Delta,y)]
\right)^{2\beta}
\varphi_T(y)\,dy.
\]
By the Cauchy--Schwarz inequality in $y$,
\begin{align}
\|F_s(B,\Delta)\|_2^2
&\leq
\left(
\int_{\R^d}
\|\nabla_bL_T(B+s\Delta,y)\|_2^2
\varphi_T(y)\,dy
\right)
I_\beta^\pi(s).
\label{eq:Fs-two-factors}
\end{align}

We first bound the first factor in \eqref{eq:Fs-two-factors}. Recall that
\[
\|\nabla_bL_T(b,y)\|_2^2
\leq
\|\sigma^{-1}\|_F^2
L_T(b,y)^2
\left(
1+\|y\|_2^2+T^2\|a(b)\|_2^2
\right).
\]
Using
\[
L_T(b,y)^2\varphi_T(y)
=
\exp\!\left(T\|a(b)\|_2^2\right)
\varphi_T(y-2Ta(b))
\]
and the corresponding Gaussian moment formula, we obtain
\[
\int_{\R^d}
\|\nabla_bL_T(B+s\Delta,y)\|_2^2
\varphi_T(y)\,dy
\leq
C
\left(
1+\|a(B+s\Delta)\|_2^2
\right)
\exp\!\left(
T\|a(B+s\Delta)\|_2^2
\right).
\]
Since
\[
a(B+s\Delta)=a(B)+s\sigma^{-1}\Delta,
\]
we have
\[
\|a(B+s\Delta)\|_2^2
\leq
2\|a(B)\|_2^2
+
2\|\sigma^{-1}\|_F^2\|\Delta\|_2^2.
\]
Consequently,
\begin{align}
&\int_{\R^d}
\|\nabla_bL_T(B+s\Delta,y)\|_2^2
\varphi_T(y)\,dy \nonumber\\
&\quad\leq
C
\left(
1+2\|a(B)\|_2^2
+2\|\sigma^{-1}\|_F^2\|\Delta\|_2^2
\right)
\exp\!\left(
2T\|a(B)\|_2^2
+2T\|\sigma^{-1}\|_F^2\|\Delta\|_2^2
\right).
\label{eq:first-factor-pre-envelope}
\end{align}

By Assumption~\ref{subGaussianFinal}, $\tau>4T\|\sigma^{-1}\|_F^2$ and $\gamma_0^2>8T\|\sigma^{-1}\|_F^2.$
We may therefore choose $\eta>0$ sufficiently small such that
\[
4(1+\eta)T\|\sigma^{-1}\|_F^2<\tau
\]
and
\[
8(1+\eta)T\|\sigma^{-1}\|_F^2<\gamma_0^2.
\]
Next, choose $\varepsilon_\Delta>0$ sufficiently small that
\[
2(1+\eta)
\left(
2T\|\sigma^{-1}\|_F^2+\varepsilon_\Delta
\right)
<\tau,
\]
and define
\[
C_1
:=
2T\|\sigma^{-1}\|_F^2+\varepsilon_\Delta.
\]
For $x\geq0$,
\[
1+2\|\sigma^{-1}\|_F^2x
\leq
K_\Delta e^{\varepsilon_\Delta x}
\]
for a finite constant $K_\Delta$. Also,
\[
1+2\|a(B)\|_2^2
+2\|\sigma^{-1}\|_F^2\|\Delta\|_2^2
\leq
\left(
1+2\|a(B)\|_2^2
\right)
\left(
1+2\|\sigma^{-1}\|_F^2\|\Delta\|_2^2
\right).
\]
It follows from \eqref{eq:first-factor-pre-envelope} that
\begin{equation}\label{eq:first-factor-envelope}
\int_{\R^d}
\|\nabla_bL_T(B+s\Delta,y)\|_2^2
\varphi_T(y)\,dy
\leq
C_0(B)
\exp\!\left(C_1\|\Delta\|_2^2\right),
\end{equation}
where
\[
C_0(B)
:=
CK_\Delta
\left(
1+2\|a(B)\|_2^2
\right)
\exp\!\left(
2T\|a(B)\|_2^2
\right).
\]

Using
\[
\|a(B)\|_2^2
\leq
2\|\sigma^{-1}\|_F^2\|B\|_2^2
+2\|m\|_2^2,
\]
we obtain, for a finite constant $K$,
\[
C_0(B)^{2(1+\eta)}
\leq
K
\left(
1+\|B\|_2^2
\right)^{2(1+\eta)}
\exp\!\left(
8(1+\eta)T
\|\sigma^{-1}\|_F^2
\|B\|_2^2
\right).
\]
The strict inequality
\[
8(1+\eta)T\|\sigma^{-1}\|_F^2<\gamma_0^2
\]
and the sub-Gaussian assumption therefore imply
\begin{equation}\label{eq:C0-moment}
\E_{\mathbb P_0}
\left[
C_0(B)^{2(1+\eta)}
\right]
<\infty.
\end{equation}

It remains to bound $I_\beta^\pi(s)$ uniformly over the couplings
considered above.

Suppose first that $\beta<0$. Define
\[
a_s(B,\Delta)
:=
a(B)+s\sigma^{-1}\Delta.
\]
By Jensen's inequality,
\begin{align*}
\E_\pi[L_T(B+s\Delta,y)]
&=
\E_\pi
\left[
\exp\!\left(
\langle a_s(B,\Delta),y\rangle
-
\frac{T}{2}\|a_s(B,\Delta)\|_2^2
\right)
\right] \\
&\geq
\exp\!\left(
\left\langle
\E_\pi[a_s(B,\Delta)],y
\right\rangle
-
\frac{T}{2}
\E_\pi
\left[
\|a_s(B,\Delta)\|_2^2
\right]
\right).
\end{align*}
Since $2\beta<0$, raising both sides to the power $2\beta$ reverses the
inequality. Integrating the resulting bound against $\varphi_T$ gives
\[
I_\beta^\pi(s)
\leq
\exp\!\left(
(2\beta^2-\beta)T
\E_\pi
\left[
\|a_s(B,\Delta)\|_2^2
\right]
\right).
\]
Moreover,
\[
\E_\pi
\left[
\|a_s(B,\Delta)\|_2^2
\right]
\leq
2\E_{\mathbb P_0}
\left[
\|a(B)\|_2^2
\right]
+
2\|\sigma^{-1}\|_F^2
\E_\pi
\left[
\|\Delta\|_2^2
\right].
\]
By Jensen's inequality and the exponential-cost bound,
\[
\exp\!\left(
\tau\E_\pi\|\Delta\|_2^2
\right)
\leq
\E_\pi
\left[
\exp\!\left(
\tau\|\Delta\|_2^2
\right)
\right]
\leq 2,
\]
and therefore
\[
\E_\pi\|\Delta\|_2^2
\leq
\frac{\log 2}{\tau}.
\]
It follows that there exists a finite constant $K_{\beta,-}$,
independent of $\pi$ and $s$, such that
\begin{equation}\label{eq:Ibeta-negative-uniform}
I_\beta^\pi(s)
\leq
K_{\beta,-},
\qquad
\beta<0.
\end{equation}

Now suppose that $\beta>0$. Choose
$r_\beta,p_\beta,q_\beta$ as in
Assumption~\ref{subGaussianFinal}. Thus,
\[
r_\beta>\max\{1,2\beta\},
\qquad
\frac{1}{p_\beta}+\frac{1}{q_\beta}=1,
\]
and
\[
2T(r_\beta^2-r_\beta)
\|\sigma^{-1}\|_F^2p_\beta
<
\gamma_0^2,
\]
while
\[
2T(r_\beta^2-r_\beta)
\|\sigma^{-1}\|_F^2q_\beta
<
\tau.
\]
By the Lyapunov and Jensen inequalities,
\[
I_\beta^\pi(s)
\leq
\left(
\E_\pi
\left[
\int_{\R^d}
L_T(B+s\Delta,y)^{r_\beta}
\varphi_T(y)\,dy
\right]
\right)^{2\beta/r_\beta}.
\]
Completing the square gives
\[
\int_{\R^d}
L_T(b,y)^{r_\beta}
\varphi_T(y)\,dy
=
\exp\!\left(
\frac{T}{2}
(r_\beta^2-r_\beta)
\|\sigma^{-1}b-m\|_2^2
\right).
\]
Since
\[
\|\sigma^{-1}(B+s\Delta)-m\|_2^2
\leq
4\|\sigma^{-1}\|_F^2
\left(
\|B\|_2^2+\|\Delta\|_2^2
\right)
+2\|m\|_2^2,
\]
{H\"older's inequality} yields
\begin{align*}
I_\beta^\pi(s)
&\leq
K_{\beta,+}
\left(
\E_{\mathbb P_0}
\left[
\exp\!\left(
2T(r_\beta^2-r_\beta)
\|\sigma^{-1}\|_F^2p_\beta
\|B\|_2^2
\right)
\right]
\right)^{\frac{2\beta}{r_\beta p_\beta}} \\
&\quad\times
\left(
\E_\pi
\left[
\exp\!\left(
2T(r_\beta^2-r_\beta)
\|\sigma^{-1}\|_F^2q_\beta
\|\Delta\|_2^2
\right)
\right]
\right)^{\frac{2\beta}{r_\beta q_\beta}},
\end{align*}
where $K_{\beta,+}<\infty$ is independent of $\pi$ and $s$.
The first expectation is finite by the sub-Gaussian assumption. The
coefficient in the second expectation is strictly smaller than $\tau$,
and hence
\[
\E_\pi
\left[
\exp\!\left(
2T(r_\beta^2-r_\beta)
\|\sigma^{-1}\|_F^2q_\beta
\|\Delta\|_2^2
\right)
\right]
\leq
\E_\pi
\left[
\exp\!\left(
\tau\|\Delta\|_2^2
\right)
\right]
\leq 2.
\]
Therefore, there exists a finite constant $\widetilde K_{\beta,+}$,
independent of $\pi$ and $s$, such that
\begin{equation}\label{eq:Ibeta-positive-uniform}
I_\beta^\pi(s)
\leq
\widetilde K_{\beta,+},
\qquad
\beta>0.
\end{equation}

Combining \eqref{eq:Ibeta-negative-uniform} and
\eqref{eq:Ibeta-positive-uniform}, there exists a finite constant
$K_\beta$, independent of $\pi$ and $s$, such that $I_\beta^\pi(s)\leq K_\beta.$
Together with \eqref{eq:Fs-two-factors} and
\eqref{eq:first-factor-envelope}, this gives
\[
\|F_s(B,\Delta)\|_2^2
\leq
K_\beta C_0(B)
\exp\!\left(
C_1\|\Delta\|_2^2
\right),
\qquad
s\in[0,1],
\]
$\pi$-almost surely.

Finally, set $C(B):=K_\beta C_0(B).$
Then \eqref{eq:C0-moment} gives $\E_{\mathbb P_0}
\left[
C(B)^{2(1+\eta)}
\right]
<\infty,$
and by construction, $2(1+\eta)C_1<\tau.$
This proves part~(2).
\end{proof}

\end{document}
